\section{Stakeholders e intereses}
La Tabla~\ref{tab:stakeholders} presenta cada stakeholder, su información de contacto y los intereses o expectativas que tiene sobre el sistema, es decir, cómo espera beneficiarse de él.

\begin{table}[htbp]
\centering
\caption{Stakeholders e intereses}
\label{tab:stakeholders}
\renewcommand{\arraystretch}{1.3}
\begin{tabularx}{\textwidth}{|>{\raggedright\arraybackslash}p{3cm}|>{\raggedright\arraybackslash}p{5.5cm}|X|}
\hline
\textbf{Rol} & \textbf{Información de contacto} & \textbf{Intereses / expectativas} \\ \hline
Autor(es) del proyecto & Andrés Henao Niño : ahenaon@javeriana.edu.co \newline Joshua Prieto Zambrano : prietojoshua@javeriana.edu.co \newline Joan Emmanuel Orduz Chia : joan\_orduz@javeriana.edu.co \newline David Estevan Rodriguez Jurado : rodriguezj.david@javeriana.edu.co & Cumplir los objetivos de negocio de la sección 1.3: interfaz unificada entre pasarelas, extensibilidad, entorno de pruebas controlado y mejores métricas de calidad de código (WMC, CBO, RFC) frente a una integración directa. Para ello se usa la arquitectura hexagonal como decisión de diseño (ADR-01, sección 13.1). Entregar funcionales el SDK, la API de Simulación y la documentación centralizada dentro del tiempo disponible. \\ \hline
Director del proyecto de grado / Asesor & Jaime Andres Pavlich Mariscal : jpavlich@javeriana.edu.co & Coherencia con el SRS y el SPMP; seguimiento del avance frente al cronograma. \\ \hline
Evaluadores académicos / Jurado de tesis & Se asigna en la sustentación & Rigor en la comparación de métricas CK; alcance delimitado con claridad (sección 7.3). \\ \hline
Desarrollador de software / Tester QA (usuarios finales) & No se define un canal de contacto formal en el alcance del proyecto. & Integración ágil (instalación y primeros pagos en <2h, sección 1.4); conmutación de pasarela por configuración (sección 1.3); pruebas sin credenciales reales mediante simulación (RF-09, RF-10); errores tipados y recuperación automática ante fallas (ADR-04, ADR-05). \\ \hline
\end{tabularx}
\end{table}

